\section{习题与解答}
\begin{exercise}\label{ex:o4twohalf}
两块理想半波片的快轴方向依次为$\alpha_1,\alpha_2$。求合成Jones矩阵和线偏振转角，说明为何要保留本征轴方向而不能只给光场乘一个共同相位。
\Diagram{fig:o4twohalfq}{两半波片的先后轴方向决定合成旋转}{\draw[axis](-1,0)--(5.5,0);\draw[curve](0,-1)--(0,1);\draw[curve](3,-1)--(3,1);\node[above]at(0,1){$H(\alpha_1)$};\node[above]at(3,1){$H(\alpha_2)$};}
\end{exercise}
\begin{solution}
删去每片的共同相位，图\ref{fig:o4twohalfq}半波片矩阵为
\[
 H(\alpha)=\begin{pmatrix}\cos2\alpha&\sin2\alpha\\
 \sin2\alpha&-\cos2\alpha\end{pmatrix}.
\]
令$c_j=\cos2\alpha_j$、$s_j=\sin2\alpha_j$，按传播顺序相乘得到
\[
 H(\alpha_2)H(\alpha_1)=
 \begin{pmatrix}c_2c_1+s_2s_1&c_2s_1-s_2c_1\\
 s_2c_1-c_2s_1&s_2s_1+c_2c_1\end{pmatrix}
 =R[2(\alpha_2-\alpha_1)].
\]
其中$R(\theta)=\bigl(\begin{smallmatrix}\cos\theta&-\sin\theta\\\sin\theta&\cos\theta\end{smallmatrix}\bigr)$。因此任意线偏振方向转过$2(\alpha_2-\alpha_1)$，椭圆的长轴同样旋转，旋向经两次翻转后恢复。
\Diagram{fig:o4twohalfa}{图取$\alpha_1=20^\circ,\alpha_2=60^\circ$，合成旋转$80^\circ$}{\draw[axis](0,0)--(3,0)node[right]{原方向};\draw[guide](0,0)--(2.53717,.92345)node[right]{$\alpha_1$};\draw[guide](0,0)--(1.5,2.59808)node[above]{$\alpha_2$};\draw[axis](0,0)--(.45149,2.56050)node[above left]{输出};\draw[->](1,0)arc(0:80:1);\node[below]at(1,-.6){$2(\alpha_2-\alpha_1)$};}
图\ref{fig:o4twohalfa}反映两次镜像的几何关系。交换波片顺序会使旋转角反号；快轴平行时合成为单位矩阵。半波片的相对延迟为$\pi$，不是对所有输入都加同一个标量相位；忽略本征轴会丢失整个旋转效应。
\end{solution}

\begin{exercise}\label{ex:o4geometricphase}
干涉仪两臂输入同一正圆偏振态$(1,i)^{\mathsf T}/\sqrt2$。一臂无偏振元件，另一臂依次通过快轴0度和$\alpha$的两块半波片。普通传播相位已补偿，两臂等强。求输出偏振、相对相位和合束强度。
\Diagram{fig:o4geometricq}{两臂最终偏振相同，波片组合仍可改变干涉相位}{\draw[->](0,1)--(5.5,1);\draw[->](0,-.5)--(5.5,-.5);\draw[curve](1.5,-1)--(1.5,0);\draw[curve](3.5,-1)--(3.5,0);\node[below]at(1.5,-1){$H(0)$};\node[below]at(3.5,-1){$H(\alpha)$};\node[right]at(5.5,.3){合束};}
\end{exercise}
\begin{solution}
由上一题，图\ref{fig:o4geometricq}测试臂的矩阵为$R(2\alpha)$。对圆Jones矢量直接作用，
\[
 R(2\alpha)\frac1{\sqrt2}\begin{pmatrix}1\\i\end{pmatrix}
 =\ee^{-2i\alpha}\frac1{\sqrt2}\begin{pmatrix}1\\i\end{pmatrix}.
\]
所以输出仍是同一正圆偏振，公共相位相对参考臂却改变$-2\alpha$。若单臂强度为$I_1$，两臂合场比例$1+\ee^{-2i\alpha}$，强度为$4I_1\cos^2\alpha$。
\Diagram{fig:o4geometrica}{两臂振幅相等，圆态的公共相位可由干涉测出}{\draw[axis](0,0)--(2.5,0)node[right]{参考相量};\draw[axis](0,0)--(1,-2.291288)node[right]{测试相量};\draw[guide](2.5,0)arc(0:-66.42182:2.5);\node at(1.4,-.5){$-2\alpha$};}
图\ref{fig:o4geometrica}的相位随波片转角改变，$\alpha=0$相长，$\alpha=\pi/2$相消。单臂的偏振椭圆保持不变，两臂的干涉强度却随这个公共相位变化。计算中删去的两片材料本身公共传播相位必须确已补偿，本题只隔离随轴角变化的部分。
\end{solution}

\begin{exercise}\label{ex:o4compensator}
完全偏振光Jones矢量为$(1,2\ee^{i\pi/3})^{\mathsf T}$。只允许转动一块四分之一波片，求可使出射成为线偏振的轴方向，并给出线偏振输出与椭圆主轴的关系。
\Diagram{fig:o4compq}{先由椭圆主轴确定补偿波片的本征轴}{\draw[axis](-1.8,0)--(2,0)node[right]{$x$};\draw[axis](0,-2.2)--(0,2.4)node[above]{$y$};\draw[curve,domain=0:360,samples=100]plot({cos(\x)},{2*cos(\x-60)});\node[right]at(1.5,1.7){待补偿};}
\end{exercise}
\begin{solution}
图\ref{fig:o4compq}有$a=1,b=2,\delta=\pi/3$。主轴角满足$\tan2\psi=-2/3$，用$\operatorname{atan2}(2,-3)$选长轴，$\psi\simeq73.155^\circ$。长短半轴平方为$A^2=(5+\sqrt{13})/2$、$B^2=(5-\sqrt{13})/2$。
\Diagram{fig:o4compa}{快慢轴对准长短轴后，90度相差被补成0或180度}{\draw[axis](-2,0)--(2.4,0)node[right]{长轴$\xi$};\draw[axis](0,-1.6)--(0,2)node[above]{短轴$\eta$};\draw[curve](0,0)ellipse(1.8 and .725);\draw[very thick](-1.8,.725)--(1.8,-.725);\draw[guide](-1.8,-.725)--(1.8,.725);}
原态$S_3>0$，在主轴基删公共相位可写$(A,iB)$。若快轴沿长轴、慢轴沿短轴，波片后为$(A,-B)$，输出线方向为$\psi-\arctan(B/A)$；若快慢互换，成为$(iA,iB)$，方向为$\psi+\arctan(B/A)$。图\ref{fig:o4compa}画出两种补偿线。因此波片快轴可沿$73.155^\circ$或其正交方向，输出方向由补偿后的两个非零分量合成，一般偏离原椭圆长轴。
\end{solution}

\begin{exercise}\label{ex:o4ellipsereconstruct}
完全偏振光经旋转线检偏器的最大、最小强度比为9:1，最大方向与$x$成30度。求椭圆长短轴比及长轴方向；仅凭这些数据能否确定旋向？
\Diagram{fig:o4reconstructq}{检偏最大方向与椭圆长轴一致}{\draw[axis](-2,0)--(2.5,0)node[right]{$x$};\draw[curve,rotate=30](0,0)ellipse(1.8 and .6);\draw[guide](-1.8,-1.039)--(1.8,1.039);\node[right]at(1.8,1.039){30度};}
\end{exercise}
\begin{solution}
图\ref{fig:o4reconstructq}完全偏振条件下，沿主轴检偏得到强度分别正比$A^2,B^2$，所以$A/B=3$，长轴为30度模180度。强度比不是半轴比9。
\Diagram{fig:o4reconstructa}{同一椭圆轨迹允许两种相反旋向}{\draw[curve](0,0)ellipse(1.6 and .6);\draw[->](1.6,0)arc(0:60:1.6 and .6);\draw[->](1.6,0)arc(0:-60:1.6 and .6);\node[right]at(2,.3){需圆偏振分析};}
图\ref{fig:o4reconstructa}两种旋向有相同线检偏曲线，需波片测$S_3$符号。若没有“完全偏振”前提，9:1还可能来自自然光与线偏振混合，不能唯一恢复椭圆半轴。
\end{solution}

\begin{exercise}\label{ex:o4halfcircular}
证明理想半波片将任意完全偏振光的椭圆旋向反转，并求其长轴方向变化。输入长轴角$\psi$，快轴角$\alpha$。
\Diagram{fig:o4halfcircq}{半波片在快慢轴上把一个分量反号}{\draw[axis](-2,0)--(2.3,0)node[right]{快轴};\draw[axis](0,-1.7)--(0,2)node[above]{慢轴};\draw[curve,rotate=25](0,0)ellipse(1.6 and .7);}
\end{exercise}
\begin{solution}
图\ref{fig:o4halfcircq}中，在波片快、慢轴基底去掉共同相位，半波片为$H=\operatorname{diag}(1,-1)$。对任意时间的实电场，它使$(E_f(t),E_s(t))$变为$(E_f(t),-E_s(t))$，即关于快轴的镜像。镜像保留长短半轴，长轴相对快轴的夹角由$\psi-\alpha$变为$-(\psi-\alpha)$，所以非圆椭圆的输出方向为
\[
 \psi_o=2\alpha-\psi\pmod\pi.
\]
\Diagram{fig:o4halfcirca}{镜像变换保留椭圆形状，反转沿轨迹的有向旋转}{\draw[axis](-2.4,0)--(2.6,0)node[right]{快轴};\draw[curve,rotate=25](0,0)ellipse(1.8 and .65);\draw[guide,rotate=-25](0,0)ellipse(1.8 and .65);\node[below]at(0,-1.8){$\psi_o=2\alpha-\psi$};}
图\ref{fig:o4halfcirca}中的轨迹形状不变，有向面积却反号。代数上，轴基底的互相关$E_fE_s^*$变为$-E_fE_s^*$，故$S_3$变号；绕传播轴旋转实坐标基底不改变$S_3$，这一结论在实验室基底同样成立。完全偏振光的椭圆率大小保持，旋向反转。

有两个退化情形须分开：线偏振的$S_3=0$，本来没有椭圆旋向，只需使用线方向公式；圆偏振满足长短半轴相同，长轴方向$\psi$没有定义，不能再用长轴角描述输出。对圆Jones矢量直接作用$H$，$(1,i)$变为$(1,-i)$，能明确判断圆旋向反转。公共相位改变时间参数的起点，不改变上述几何结论。
\end{solution}

\begin{exercise}\label{ex:o4twoquarter}
两块相同四分之一波片快轴平行或正交。分别求合成Jones矩阵，并说明两片都存在时偏振是否一定改变。
\Diagram{fig:o4twoquarterq}{两片波片的相对轴角决定相位差相加或相消}{\draw[curve](0,-1)--(0,1);\draw[curve](3,-1)--(3,1);\draw[axis](-1,0)--(4.5,0);\node[above]at(0,1){$Q_1$};\node[above]at(3,1){$Q_2$};}
\end{exercise}
\begin{solution}
按图\ref{fig:o4twoquarterq}，快轴平行时两矩阵均为$\operatorname{diag}(1,i)$，乘积$\operatorname{diag}(1,-1)$，等效半波片。快轴正交时第二片在同一基底为$\operatorname{diag}(i,1)$，乘积$iI$。
\Diagram{fig:o4twoquartera}{正交轴相消后只剩共同相位，所有输入偏振态保持}{\node at(0,.5){平行：$Q^2=H$};\node at(0,-.6){正交：$Q_\perp Q=iI$};\draw[->](2,.5)--(3.3,.5);\draw[->](2,-.6)--(3.3,-.6);\node[right]at(3.5,.5){相差累加};\node[right]at(3.5,-.6){偏振不变};}
图\ref{fig:o4twoquartera}的正交轴组合在所用约化矩阵中为$iI$，所有输入偏振态保持。若恢复每片删去的快轴传播相位$\phi_{f1},\phi_{f2}$，完整矩阵为$\ee^{i(\phi_{f1}+\phi_{f2})}iI$；与参考臂干涉时须保留这整个共同相位，不能仅由$iI$断言实物插入后的总相移恰为$\pi/2$。平行轴组合虽等效半波片，对沿其本征轴的输入仍只改变共同相位。
\end{solution}

\begin{exercise}\label{ex:o4criticaldecay}
全反射入射角为$i=i_c+\epsilon$，$0<\epsilon\ll1$，$i_c=\arcsin(n_2/n_1)$。求倏逝振幅衰减长度随$\epsilon$的领先阶，说明临界角附近的尺度。
\Diagram{fig:o4criticalq}{刚超过临界角，低折射率侧场衰减很慢}{\draw[curve](-2.3,0)--(2.3,0);\draw[guide](0,-1.5)--(0,2);\draw[->](-1.6,-1.3)--(0,0);\draw[->](0,0)--(1.6,-1.3);\node[right]at(1.2,1){长衰减尾部};}
\end{exercise}
\begin{solution}
图\ref{fig:o4criticalq}的$\kappa=k_0\sqrt{n_1^2\sin^2i-n_2^2}$。记$F(i)=n_1^2\sin^2i-n_2^2$，对平方和三角函数逐层求导有$F'(i)=2n_1^2\sin i\cos i$。利用$F(i_c)=0$、$\sin i_c=n_2/n_1$，得到
\begin{align*}
 F(i_c+\epsilon)&=F(i_c)+F'(i_c)\epsilon+O(\epsilon^2)\\
 &=2n_2\sqrt{n_1^2-n_2^2}\,\epsilon+O(\epsilon^2).
\end{align*}
将这一领先项代入$\ell=1/(k_0\sqrt F)$，得
\[
 \ell=\frac1\kappa\simeq
 \frac{\lambda_0}{2\pi\sqrt{2n_2\sqrt{n_1^2-n_2^2}\,\epsilon}}.
\]
\Diagram{fig:o4criticala}{衰减长度在临界角右侧按$\epsilon^{-1/2}$增大}{\draw[axis](0,0)--(4.5,0)node[right]{$\epsilon$};\draw[axis](0,0)--(0,3.5)node[above]{$\ell$};\draw[curve,domain=.2:4,samples=70]plot(\x,{1.4/sqrt(\x)});}
图\ref{fig:o4criticala}的发散说明接近临界角时必须考虑有限光束角谱和第二界面距离。若光束有角度展宽，一部分成分可能低于临界角，单一平面波的无穷衰减长度不再完整描述实际束形。
\end{solution}

\begin{exercise}\label{ex:o4magneticinterface}
正常入射到无损磁性介质界面，介质波阻抗分别为$Z_1=\sqrt{\mu_1/\varepsilon_1}$、$Z_2=\sqrt{\mu_2/\varepsilon_2}$。推导共同切向电场基底的$r,t,R,T$，说明何时可只用折射率。
\Diagram{fig:o4magneticq}{正常入射时电场与磁场的比例由波阻抗决定}{\draw[curve](2,-1.3)--(2,1.3);\draw[->](0,.5)--(1.9,.5);\draw[->](1.9,-.4)--(0,-.4);\draw[->](2.1,.5)--(4.5,.5);\node[above]at(.8,1){$Z_1$};\node[above]at(3.3,1){$Z_2$};}
\end{exercise}
\begin{solution}
图\ref{fig:o4magneticq}中，统一取切向电场沿$x$，入射波沿$+z$、反射波沿$-z$。由$\vct H=\hat{\vct k}\times\vct E/Z$，反射磁场相对其电场多一个方向反号。将入射电场幅度取1，边界条件为
\[
 1+r=t,\qquad \frac{1-r}{Z_1}=\frac t{Z_2}.
\]
用$t=1+r$代入，得到$Z_2-Z_2r=Z_1+Z_1r$，因此$r=(Z_2-Z_1)/(Z_2+Z_1)$、$t=2Z_2/(Z_1+Z_2)$。
\Diagram{fig:o4magnetica}{阻抗匹配可无反射，即使相折射率不同}{\draw[curve](0,-1.5)--(0,1.7);\node[above]at(-1.7,1.5){$Z_1$};\node[above]at(1.8,1.5){$Z_2=Z_1$};\draw[axis](-3,0)--(3.2,0)node[right]{$z$};\draw[->,thick](-2.3,0)--(-2.3,1)node[above]{$E_i$};\draw[->,thick](1.5,0)--(1.5,1)node[above]{$E_t=E_i$};\node[below]at(-1.6,-.6){$E_r=0$};\node[below]at(1.5,-.6){$T=1$};}
无损介质$Z_1,Z_2>0$，平面波平均能流大小为$|E|^2/(2Z)$。图\ref{fig:o4magnetica}中反、入射处于同一介质，透射处于介质2，所以
\[
 R=|r|^2,\qquad T=\frac{Z_1}{Z_2}|t|^2
 =\frac{4Z_1Z_2}{(Z_1+Z_2)^2}.
\]
两者相加的分子为$(Z_2-Z_1)^2+4Z_1Z_2=(Z_1+Z_2)^2$，故$R+T=1$。当$Z_1=Z_2$时无反射，即使两介质的折射率不同也成立，因为阻抗依赖$\mu/\varepsilon$，折射率依赖$\mu\varepsilon$。两介质均非磁性、磁导率相同时，$Z\propto1/n$，才可把阻抗比换成反折射率比，恢复通常正常入射公式。
\end{solution}

\begin{exercise}\label{ex:o4uncertaintystate}
两束互不相干光分别为正、负圆偏振，强度为$3I$和$I$。求总斯托克斯量和偏振度；旋转线检偏器能否观察到强度变化？
\Diagram{fig:o4mixtureq}{两相反圆偏振分量不相干地叠加}{\begin{scope}[xshift=-2cm]\draw[curve](0,0)circle(.8);\draw[->](.8,0)arc(0:70:.8);\node[below]at(0,-1.1){$3I$};\end{scope}\begin{scope}[xshift=2cm]\draw[curve](0,0)circle(.8);\draw[->](.8,0)arc(0:-70:.8);\node[below]at(0,-1.1){$I$};\end{scope}}
\end{exercise}
\begin{solution}
图\ref{fig:o4mixtureq}两圆态各有$S_1=S_2=0$，$S_3$分别等于其强度的正、负值。直接加斯托克斯量，得$S_0=4I,S_1=S_2=0,S_3=2I$，偏振度$1/2$。
\Diagram{fig:o4mixturea}{线检偏透过恒为$2I$，波片分析后才显出圆偏振不平衡}{\draw[axis](0,0)--(4.6,0)node[right]{检偏角};\draw[axis](0,0)--(0,2.8)node[above]{透射强度};\draw[very thick](0,1.7)--(4,1.7);\node[left]at(0,1.7){$2I$};}
线检偏强度恒为$S_0/2=2I$，如图\ref{fig:o4mixturea}。虽然它的旋转曲线与自然光相同，光束仍有一半的圆偏振成分；加四分之一波片后两线方向透过强度成为$3I$与$I$，即可识别。独立光束按斯托克斯量相加；Jones矢量直接相加只适用于有确定相位关系的场。
\end{solution}

\begin{exercise}\label{ex:o4stack}
自然光从空气入射，穿过一叠相互平行的无损玻璃片后回到空气，每个界面均处于相应方向的布儒斯特角。设共经历$M$个界面（$M$为偶数），忽略反射光重新进入透射主束。若每界面$s$偏振反射率为$R_s$，求直接透射主束光强和线偏振度；另说明若在玻璃内部截取主束，光强应怎样换算。
\Diagram{fig:o4stackq}{玻片堆逐渐削弱$s$分量，保留$p$分量}{\foreach\x in{0,1,2,3}{\draw[curve](\x,-.9)--({\x+.5},.9);}\draw[axis](-1,0)--(4.5,0)node[right]{主束};\node[above]at(1.8,1.3){每面为布儒斯特角};}
\end{exercise}
\begin{solution}
图\ref{fig:o4stackq}的自然光初始$p,s$各占$I_0/2$。每界面的法向能流透射率为$T_p=1$、$T_s=1-R_s$，而沿光束方向定义的强度还须换算截面。若第$j$面前、后的光线与共同法线分别成$\theta_{j-1},\theta_j$，有
\[
 I_{\sigma,j}\cos\theta_j
 =T_{\sigma,j}I_{\sigma,j-1}\cos\theta_{j-1},
 \qquad \sigma=p,s.
\]
逐面相乘，所有中间余弦因子抵消。最终又回到空气，$\theta_M=\theta_0$，所以
\[
 I_p=\frac{I_0}{2},\qquad
 I_s=\frac{I_0}{2}(1-R_s)^M.
\]
\Diagram{fig:o4stacka}{两正交强度份额决定偏振度}{\draw[curve](0,0)rectangle(4,.6);\node at(2,.3){$I_p=I_0/2$};\draw[curve](0,-1)rectangle(1.5,-.4);\node[right]at(1.7,-.7){$I_s=(I_0/2)(1-R_s)^M$};}
总强度$I=I_0[1+(1-R_s)^M]/2$，线偏振度$\mathcal P=[1-(1-R_s)^M]/[1+(1-R_s)^M]$。图\ref{fig:o4stacka}在$0<R_s<1$、$M\to\infty$时趋于完全$p$偏振、主束强度$I_0/2$。

若在玻璃内的第$j$面后截取主束，上述两个强度须共同乘$\cos\theta_0/\cos\theta_j$，并将幂次$M$改为$j$；偏振度中共同因子抵消。这样便保留了法向通量与光束强度的区别。若多次反射旁束也进入探测孔径，须重新求和；有吸收或表面误差时，$p$分量也会损失。
\end{solution}

\begin{exercise}\label{ex:o4partialmeasure}
未知光由自然光和某方向线偏振光组成。旋转检偏器得到最大强度8、最小强度2，单位相同。求入射总强度、自然成分和线偏振成分。若去掉“线偏振”这一前提，还能唯一确定偏振度吗？
\Diagram{fig:o4partialq}{检偏器扫描得到非零的最低强度}{\draw[axis](0,0)--(5.2,0)node[right]{角度};\draw[axis](0,0)--(0,3)node[above]{透射强度};\draw[curve,domain=0:180,samples=80]plot({\x/40},{(2+6*cos(\x)^2)/3});\node[left]at(0,2.667){$8$};\node[left]at(0,.667){$2$};}
\end{exercise}
\begin{solution}
按图\ref{fig:o4partialq}，$I_u/2=2$，所以$I_u=4$；$I_p+I_u/2=8$给$I_p=6$。入射总强度$I_0=10$，偏振度$6/10=.6$。
\Diagram{fig:o4partiala}{线检偏扫描测得$S_1,S_2$，不能单独得到$S_3$}{\draw[curve](0,0)rectangle(2.5,1);\node at(1.25,.5){线检偏器};\draw[->](2.6,.5)--(3.5,.5);\node[right]at(3.5,.5){$S_0,S_1,S_2$};\node[below]at(2,-.6){圆偏振信息需波片};}
若完全偏振分量允许为椭圆或圆偏振，图\ref{fig:o4partiala}说明圆偏振信息$S_3$未测得。单次扫描只能给线性调制度.6，即完整偏振度的下限；它可以介于.6与1之间。附加适当四分之一波片测圆偏振分析强度才能完成判定。
\end{solution}

\begin{exercise}\label{ex:o4plateorder}
线偏振沿$x$，先过快轴45度的四分之一波片，再过快轴0度的半波片。求输出Jones矢量，并比较交换两片顺序的结果。
\Diagram{fig:o4orderq}{两块转角不同的波片，矩阵按传播次序从右向左相乘}{\draw[axis](-1,0)--(5.5,0);\draw[curve](0,-1)--(0,1);\draw[curve](2.8,-1)--(2.8,1);\node[above]at(0,1){$Q(45^\circ)$};\node[above]at(2.8,1){$H(0)$};\node[left]at(-1,.5){$x$线偏振};}
\end{exercise}
\begin{solution}
由旋转矩阵，$Q(45^\circ)=\tfrac12\bigl(\begin{smallmatrix}1+i&1-i\\1-i&1+i\end{smallmatrix}\bigr)$，$H(0)=\operatorname{diag}(1,-1)$。图\ref{fig:o4orderq}的输出为$\tfrac12(1+i,-1+i)^{\mathsf T}$，删去共同相位可写$(1,i)^{\mathsf T}/\sqrt2$。
\Diagram{fig:o4ordera}{交换顺序使最终圆偏振的旋向相反}{\begin{scope}[xshift=-2cm]\draw[curve](0,0)circle(.9);\draw[->](.9,0)arc(0:70:.9);\node[below]at(0,-1.2){先$Q$后$H$};\end{scope}\begin{scope}[xshift=2cm]\draw[curve](0,0)circle(.9);\draw[->](.9,0)arc(0:-70:.9);\node[below]at(0,-1.2){先$H$后$Q$};\end{scope}}
交换后$H(0)$先作用于$(1,0)$不改偏振，输出为$\tfrac12(1+i,1-i)^{\mathsf T}$，删共同相位为$(1,-i)^{\mathsf T}/\sqrt2$。图\ref{fig:o4ordera}的两箭头相反；不对易来自两波片的本征轴不同，而不是共同相位差。
\end{solution}

\begin{exercise}\label{ex:o4normalenergy}
空气向折射率$n=1.5$玻璃正常入射。求振幅反射、透射系数以及能量反射、透射率。再解释为何透射场振幅小于入射场并不表示透过能量只有振幅比的平方。
\Diagram{fig:o4normalq}{正常入射界面，采用共同切向电场基底}{\draw[curve](2,-1.4)--(2,1.4);\draw[->](0,.6)--(1.9,.6);\draw[->](1.9,0)--(0,0);\draw[->](2.1,.6)--(4.5,.6);\node[above]at(.8,1){$n_1=1$};\node[above]at(3.3,1){$n_2=1.5$};}
\end{exercise}
\begin{solution}
图\ref{fig:o4normalq}统一切向场有$r=(1-1.5)/(1+1.5)=-.2$、$t=2/(1+1.5)=.8$。所以$R=.04$，$T=(n_2/n_1)t^2=1.5\times.64=.96$。
\Diagram{fig:o4normala}{能量份额为4\%反射、96\%透射}{\draw[curve](0,0)rectangle(5,1);\draw[curve](.2,0)--(.2,1);\node[above]at(.1,1){4\%};\node at(2.6,.5){96\%};\node[below]at(2.5,0){$R+T=1$};}
图\ref{fig:o4normala}满足能量守恒。平面波光强为$n\varepsilon_0cE_0^2/2$，介质2具有不同波导纳，因此$t^2=.64$只是电场振幅平方比。若使用局部$p$基，振幅$r_p=+.2$，能量反射率仍相同，必须同时注明基向量改变。
\end{solution}

\begin{exercise}\label{ex:o4brewsternumerical}
自然光从空气射向$n=1.5$玻璃，入射角为布儒斯特角。求该角、折射角、反射光强相对入射光强的比例及反射偏振方向。
\Diagram{fig:o4brewsterq}{布儒斯特条件使反射与透射光线垂直}{\draw[curve](-2.5,0)--(2.5,0);\draw[guide](0,-1.7)--(0,2);\draw[->](-2,-1.333)--(0,0);\draw[->](0,0)--(2,-1.333);\draw[->](0,0)--(1.2,1.8);\node[right]at(2,-1.333){反射};\node[right]at(1.2,1.8){透射};}
\end{exercise}
\begin{solution}
图\ref{fig:o4brewsterq}满足$\tan i_B=1.5$，所以$i_B\simeq56.31^\circ$、$r_B\simeq33.69^\circ$。$p$反射为零，$s$振幅系数代入$\cos i_B=1/\sqrt{1+n^2}$、$\cos r_B=n/\sqrt{1+n^2}$得$r_s=(1-n^2)/(1+n^2)=-5/13$。
\Diagram{fig:o4brewstera}{自然光的两个偏振份额中仅$s$有反射}{\draw[curve](0,0)rectangle(2.8,.7);\node at(1.4,.35){$I_0/2$的$s$};\draw[->](2.9,.35)--(4,.35);\node[right]at(4,.35){乘$25/169$};\node[below]at(2,-.7){$p$反射为0};}
自然光的$s$份额只有$I_0/2$，故总反射率为$(1/2)(25/169)=25/338\simeq.07396$，如图\ref{fig:o4brewstera}。反射光完全为$s$偏振，电场垂直入射面；不能把$R_s$直接当自然光的总反射率。
\end{solution}

\begin{exercise}\label{ex:o4retardthickness}
某波片在目标波长$\lambda_0$处双折射差为$\Delta n$。求零级四分之一波片厚度及厚度误差$\Delta h$引起的相位误差。若忽略材料色散，偏离波长时会怎样？
\Diagram{fig:o4thicknessq}{波片中快慢分量走过相同厚度却积累不同相位}{\draw[curve](1,-1)rectangle(3,1);\draw[->](0,.4)--(4.5,.4)node[right]{快};\draw[->](0,-.4)--(4.5,-.4)node[right]{慢};\draw[<->](1,-1.4)--(3,-1.4)node[midway,below]{$h$};}
\end{exercise}
\begin{solution}
图\ref{fig:o4thicknessq}中令$\Delta n=n_s-n_f>0$，慢轴相对快轴的相位延迟为$\delta=2\pi\Delta n h/\lambda_0$。零级四分之一波片取最小正延迟$\pi/2$，故$h_0=\lambda_0/(4\Delta n)$。固定波长、材料时厚度误差给
\[
 \Delta\delta=\frac{2\pi\Delta n}{\lambda_0}\Delta h
 =\frac\pi2\frac{\Delta h}{h_0}.
\]
若快慢轴命名相反，延迟符号随之改变，最小厚度应使用$|\Delta n|$。
\Diagram{fig:o4thicknessa}{延迟偏离90度后，45度线偏振不再变成严格圆偏振}{\draw[axis](-1.5,0)--(1.7,0);\draw[axis](0,-1.5)--(0,1.7);\draw[guide](0,0)circle(1);\draw[curve,rotate=45](0,0)ellipse(1.2 and .75);\node[right]at(1.8,.7){圆变椭圆};}
图\ref{fig:o4thicknessa}中，45度线偏振入射后的Jones矢量正比$(1,\ee^{i\delta})$。若$\delta=\pi/2+\epsilon$，其归一化$S_3/S_0=\sin\delta\simeq1-\epsilon^2/2$，$S_2/S_0=\cos\delta\simeq-\epsilon$，不再严格为圆偏振。

固定厚度且忽略色散时，$\delta(\lambda)=(\pi/2)\lambda_0/\lambda$，从$1/(\lambda_0+\Delta\lambda)$一阶展开得$\Delta\delta\simeq-(\pi/2)\Delta\lambda/\lambda_0$。保留材料色散则应微分整个比值：
\begin{align*}
 \frac{\dd\delta}{\dd\lambda}
 &=2\pi h\left[\frac1\lambda\frac{\dd\Delta n}{\dd\lambda}-\frac{\Delta n}{\lambda^2}\right]\\
 &=\delta\left[\frac1{\Delta n}\frac{\dd\Delta n}{\dd\lambda}-\frac1\lambda\right].
\end{align*}
多级波片的设计延迟为$\pi/2+2m\pi$，其相对厚度误差或相对波长偏移对应更大的绝对相位误差。容差由完整延迟的导数确定；先模去$2\pi$会丢失多级波片的灵敏度。
\end{solution}
